\documentclass[10pt,a4paper]{amsart}
\usepackage[margin=19mm]{geometry}
\usepackage{amsmath,amssymb,mathtools}
\usepackage[scr=rsfs,cal=euler]{mathalfa}
\usepackage{xfrac}
\usepackage[unicode,hidelinks,bookmarks=false]{hyperref}
\newcommand{\ie}{i.e.\ }
\newcommand{\cf}{cf.\ }
\newcommand{\resp}{respectively\ }
\newcommand{\Subsection}{\subsection}
\providecommand{\nobreakdash}{\nobreak\hbox{-}}
\setlength{\emergencystretch}{2em}
\multlinegap0pt
\usepackage{float}
\usepackage{amssymb}
\newtheorem{assumption}{Assumption}
\newtheorem{lemma}{Lemma}
\title{\bf Convergence in probability of numerical solutions of a highly nonlinear delayed stochastic interest rate model}
\author{Emmanuel Coffie}
\date{Source: arXiv:2510.04092v3 (CC BY 4.0)}
\begin{document}
\maketitle

\noindent\emph{Source and licence.} \bf Convergence in probability of numerical solutions of a highly nonlinear delayed stochastic interest rate model. Version arXiv:2510.04092v3 (CC BY 4.0), distributed by arXiv under the Creative Commons Attribution 4.0 International licence, http://creativecommons.org/licenses/by/4.0/, as recorded in the arXiv licence metadata for this exact version and shown on the abstract page https://arxiv.org/abs/2510.04092v3. Full licence notice: \texttt{licenses/arxiv-2510.04092-CC-BY-4.0.txt}.\par\smallskip

\noindent\textit{Editorial scope.} Counted unit: the article's lemma sub3 (source0.tex lines 409--419). The article's complete original proof of it is delivered with it (source0.tex lines 421--481, one \texttt{proof} environment of 61 lines). No mathematics is written, repaired or re-derived: the statement and the proof are the article's own lines, in the article's own notation, and no retained line is substituted or re-flowed.  Retained uncounted as the source-local context the proof needs: lines 90--92; lines 108--114; lines 115--120; lines 124--132; lines 206--208; lines 344--346; lines 367--369; lines 383--385; lines 389--395.  Nothing was omitted from any retained span, so the package carries no omission record.  No retained line contains a citation, so the wrapper carries no bibliography and no bibliography entry.  No retained line contains a figure environment, an image-inclusion command or drawing code, so no image is delivered and the package declares no asset dependency.  Editorial formatting only, in addition: an AMS \texttt{amsart} preamble that restates the article's own declarations and macros used by the retained lines, namely the environments \texttt{assumption}, \texttt{lemma} on separate counters, together with the standard packages those lines need.  The retained environments print in this document as Assumption 1, Lemma 1 in order of appearance, whereas the article prints the same items with the numbering of its own full document.  The complete native source of the article as distributed in the arXiv e-print for this exact version is bundled unchanged as \texttt{sources/186413/source0.tex}.
\begin{equation}\label{cir4}
dx(t)=\alpha(\mu-x(t)^{\gamma})dt+\sigma x(t-\tau)^{r}x(t)^{\theta}dB(t),
\end{equation}

\begin{assumption} \label{sec2:assump:1} 
The parameters of \textup{SDDE} \eqref{cir4} satisfy 
\begin{equation}
1+ \gamma > 2(r+\theta),
\end{equation}
where  $\theta, r>0$ and $\gamma>1$.
\end{assumption}

\begin{assumption}\label{sec2:assump:2} 
There exist constants $D>0$ and $\ell\in (0,1]$ such that for all $ -\tau\leq u \leq t \leq 0$, the initial function $\xi$ satisfies
\begin{equation}\label{eq:15}
  |\xi(t)-\xi(u)|\le D|t-u|^{\ell}.
\end{equation}
\end{assumption}

\begin{lemma}\label{lemma1}
Let Assumption \ref{sec2:assump:1} hold. For any $R>0$, there exists a constant $G_R>0$ such that the coefficients of the \textup{SDDE} \eqref{cir4} satisfy
\begin{align}
| f(x)-f(\bar{x}) | + | g(x,y)-g(\bar{x},\bar{y}) |
\le G_R \big( |x-\bar{x}| + |y-\bar{y}| \big)
\end{align}
for all $x,y,\bar{x},\bar{y}\in \mathbb{R}$ with 
$|x|\vee |\bar{x}|\vee |y|\vee |\bar{y}| \le R$.
\end{lemma}

\begin{equation}\label{sec2:eq1}
V(x)=x^{\beta}-1-\beta\log(x), 
\end{equation}

\begin{equation}\label{sec4:eq:1}
\sup_{\vert x\vert\vee\vert y\vert \le u}\Big(|f(x)|\vee g(x,y)\Big)\le z(u), 
\end{equation}

\begin{align}\label{sec4:eq:3}
|f_{\Delta}(x)|\vee g_{\Delta}(x,y)\le z(z^{-1}(\psi(\Delta)))=\psi(\Delta),
\end{align}

\begin{equation}\label{eq:30}
dx_{\Delta}(t)= f_{\Delta}(\bar{x}_{\Delta}(t))dt+g_{\Delta}(\bar{x}_{\Delta}(t),\bar{x}_{\Delta}(t-\tau))dB(t).
\end{equation}

\begin{lemma}\label{sec5:eq:L1}
For any fixed $\Delta\in (0,\Delta^*]$ and $p\ge 2$, we have 
\begin{equation}
 \mathbb{E}\vert x_{\Delta}(t)-\bar{x}_{\Delta}(t)\vert^{p}\le c_p\Delta^{\frac{p}{2}}(\psi(\Delta))^p,
\end{equation}
for all $t\geq 0$, where $c_p$ is a generic constant that is dependent only on $p$.
\end{lemma}

\begin{lemma}\label{sub3}
Let Assumptions \ref{sec2:assump:1} and \ref{sec2:assump:2}  hold and $T>0$ be fixed. For any sufficiently large integer $k>0$, define the stopping time by
\begin{equation}\label{sec5:eq:tau}
\eta_{\Delta}=\inf\{t\in [0,T]:x_{\Delta}(t)\notin [1/k,k]\}.
\end{equation}
Then for any fixed $\Delta\in(0,\Delta^*]$, we have
\begin{equation}\label{probability2}
 \mathbb{P}(\eta_{\Delta}\leq T)\leq \frac{V(\xi(0))+K_1T+K_3D\Delta^{\ell}+c_p(K_2+K_3)\Delta^{1/2}\psi(\Delta)T}{V(1/k)\wedge V(k)},
\end{equation}
where  $K_1$, $K_2$ and $K_3$ are generic constants and $V$ is defined in \eqref{sec2:eq1}.
\end{lemma}

\begin{proof}
For $t_1\in[0,T]$, we apply the It\^{o} formula to \eqref{eq:30} to compute
\begin{align*}
&\mathbb{E}(V(x_{\Delta}(t\wedge \eta_{\Delta})))-V(\xi(0))\\
&=\mathbb{E}\int_{0}^{t_1\wedge \eta_{\Delta}}\Big(V_x(x_{\Delta}(s))f_{\Delta}(\bar{x}_{\Delta}(s))+\frac{1}{2}V_{xx}(x_{\Delta}(s))g_{\Delta}(\bar{x}_{\Delta}(s),\bar{x}_{\Delta}(s-\tau))^2\Big)ds\\
&\le\mathbb{E}\int_{0}^{\eta_{\Delta}\wedge t_1}\Big(V_x(x_{\Delta}(s))f_{\Delta}(x_{\Delta}(s))+\frac{1}{2}V_{xx}(x_{\Delta}(s))g_{\Delta}(x_{\Delta}(s),x_{\Delta}(s-\tau))^2\Big)ds\\
&+\mathbb{E}\int_{0}^{\eta_{\Delta}\wedge t_1}V_x(x_{\Delta}(s))\Big(f_{\Delta}(\bar{x}_{\Delta}(s))-f_{\Delta}(x_{\Delta}(s))\Big)ds\\
&+\mathbb{E}\int_{0}^{\eta_{\Delta}\wedge t_1}\frac{1}{2}V_{xx}(x_{\Delta}(s))\Big(g_{\Delta}(\bar{x}_{\Delta}(s),\bar{x}_{\Delta}(s-\tau))^2-g_{\Delta}(x_{\Delta}(s),x_{\Delta}(s-\tau))^2\Big)ds.
\end{align*}
By recalling the definition of the truncated functions in \eqref{sec4:eq:1}, we note that
\begin{equation}\label{sec5:eq:recall0}
f_{\Delta}(\cdot)=f(\cdot)\text{ and } g_{\Delta}(\cdot,\cdot)=g(\cdot,\cdot).
\end{equation}
Also, for  $s\in[0,\eta_{\Delta}\wedge t_1]$ with $x_{\Delta}(s),\bar{x}_{\Delta}(s),x_{\Delta}(s-\tau),\bar{x}_{\Delta}(s-\tau)\in[1/k,k]$, we observe that
\begin{equation}\label{sec5:eq:recall1}
g(\bar{x}_{\Delta}(s),\bar{x}_{\Delta}(s-\tau))\vee g(x_{\Delta}(s),x_{\Delta}(s-\tau))\le z(k).
\end{equation}
By Assumption \ref{sec2:assump:2}, we have 
\begin{align*}
&\mathbb{E}(V(x_{\Delta}(t\wedge \eta_{\Delta})))-V(\xi(0))\le K_1T+\mathbb{E}\int_{0}^{\eta_{\Delta}\wedge t_1}V_x(x_{\Delta}(s))\vert f(\bar{x}_{\Delta}(s))-f(x_{\Delta}(s))\vert ds\\
&+\mathbb{E}\int_{0}^{\eta_{\Delta}\wedge t_1}\frac{1}{2}V_{xx}(x_{\Delta}(s))\Big(g(\bar{x}_{\Delta}(s),\bar{x}_{\Delta}(s-\tau))^2-g(x_{\Delta}(s),x_{\Delta}(s-\tau))^2\Big)ds,
\end{align*}
where $LV(x_{\Delta}(s),x_{\Delta}(s-\tau))\le K_1$ for $s\in[0,t\wedge \eta_{\Delta}]$. By  an elementary inequality, we have
\begin{align*}
&\mathbb{E}(V(x_{\Delta}(t\wedge \eta_{\Delta})))-V(\xi(0))\le K_1T+\mathbb{E}\int_{0}^{\eta_{\Delta}\wedge t_1}V_x(x_{\Delta}(s))\vert f(\bar{x}_{\Delta}(s))-f(x_{\Delta}(s))\vert ds\\
&+\mathbb{E}\int_{0}^{\eta_{\Delta}\wedge t_1}\frac{1}{2}V_{xx}(x_{\Delta}(s))\Big(\vert g(\bar{x}_{\Delta}(s),\bar{x}_{\Delta}(s-\tau))-g(x_{\Delta}(s),x_{\Delta}(s-\tau))\vert\\
&\times \vert g(\bar{x}_{\Delta}(s),\bar{x}_{\Delta}(s-\tau))+g(x_{\Delta}(s),x_{\Delta}(s-\tau))\vert \Big) ds.
\end{align*}
By \eqref{sec4:eq:3}, \eqref{sec5:eq:recall0}, \eqref{sec5:eq:recall1} and Lemma \ref{lemma1}, we now have
\begin{align*}
\mathbb{E}(V(x_{\Delta}(t\wedge \eta_{\Delta})))&\le V(\xi(0))+K_1T+\mathbb{E}\int_{0}^{t_1\wedge\eta_{\Delta}}G_kV_x(x_{\Delta}(s))\vert \bar{x}_{\Delta}(s)-x_{\Delta}(s)\vert ds\\
&+\mathbb{E}\int_{0}^{\eta_{\Delta}\wedge t_1}(z(k))^2G_kV_{xx}(x_{\Delta}(s))\vert \bar{x}_{\Delta}(s)-x_{\Delta}(s)\vert ds\\
&+\mathbb{E}\int_{0}^{\eta_{\Delta}\wedge t_1}(z(k))^2G_kV_{xx}(x_{\Delta}(s))\vert \bar{x}_{\Delta}(s-\tau)-x_{\Delta}(s-\tau)\vert ds\\
&\le V(\xi(0))+K_1T+K_2\mathbb{E}\int_0^{\eta_{\Delta}\wedge t_1} \vert \bar{x}_{\Delta}(s)-x_{\Delta}(s)\vert ds\\
&+K_3\mathbb{E}\int_{0}^{\eta_{\Delta}\wedge t_1}\vert \bar{x}_{\Delta}(s-\tau)-x_{\Delta}(s-\tau)\vert ds,
\end{align*}
where
\begin{equation*}
K_2=\max_{1/k\le x\le k}\Big(G_kV_x(x)+(z(k))^2G_kV_{xx}(x)\Big)
\end{equation*}
and
\begin{equation*}
K_3=\max_{1/k\le x\le k}\Big((z(k))^2G_kV_{xx}(x)\Big).
\end{equation*}
So by Assumption \ref{sec2:assump:1} and Lemma \ref{sec5:eq:L1}, we get
\begin{align*}
\mathbb{E}(V(x_{\Delta}(t\wedge \eta_{\Delta})))&\le V(\xi(0))+K_1T+K_2\mathbb{E}\int_0^{\eta_{\Delta}\wedge t_1} \vert \bar{x}_{\Delta}(s)-x_{\Delta}(s)\vert ds\\
&+K_3\mathbb{E}\int_{-\tau}^{\eta_{\Delta}\wedge t_1}\vert \bar{x}_{\Delta}(s)-x_{\Delta}(s)\vert ds\\
&\le V(\xi(0))+K_1T+
K_2\mathbb{E}\int_{0}^{\eta_{\Delta}\wedge t_1} \vert \bar{x}_{\Delta}(s)-x_{\Delta}(s)\vert ds\\
&+K_3\mathbb{E}\int_{-\tau}^{0}|\xi([s/\Delta]\Delta)-\xi(s)|ds+K_3\mathbb{E}\int_{0}^{\eta_{\Delta}\wedge t_1}\vert \bar{x}_{\Delta}(s)-x_{\Delta}(s)\vert ds\\
&\le  V(\xi(0))+K_1T+K_3\int_{-\tau}^0\mathbb{E}|\xi([s/\Delta]\Delta)-\xi(s)|ds\\
&+(K_2+K_3)\int_{0}^T(\mathbb{E}\vert \bar{x}_{\Delta}(s)-x_{\Delta}(s)\vert^p)^{1/p}ds\\
&\le V(\xi(0))+K_1T+K_3D\Delta^{\ell}+c_p(K_2+K_3)\Delta^{1/2}\psi(\Delta)T.
\end{align*}
This implies that 
\begin{equation*}
 \mathbb{P}(\eta_{\Delta}\leq T)\leq \frac{V(\xi(0))+K_1T+K_3D\Delta^{\ell}+c_p(K_2+K_3)\Delta^{1/2}\psi(\Delta)T}{V(1/k)\wedge V(k)},
\end{equation*}
as required.
\end{proof}

\end{document}
